Генерация юмора является удобным способом изучения ограничений современных методов постобучения языковых моделей. В таких областях, как математика или программирование, сгенерированный ответ часто можно проверить внешней процедурой: решение либо совпадает с правильным ответом, либо программа проходит тесты или не проходит их. Юмор устроен иначе. Шутка может быть грамматически правильной, связной и соответствующей теме, но при этом не производить нужного комического эффекта.

При этом проблема не сводится только к субъективности юмора. Один и тот же запрос может допускать множество разных удачных шуток, тогда как наблюдаемый набор данных обычно содержит лишь небольшую их часть. Поэтому юмор является не только сложной задачей генерации, но и полезным примером для изучения рассуждения и обучения в неверифицируемых доменах \citep{Loakman2025,lemmens2026computational}. Эта сложность стала особенно важной после появления больших языковых моделей.

Ранние системы вычислительного юмора были узкими, но интерпретируемыми: они опирались на шаблоны, лексические ресурсы или конкретные механизмы каламбуров \citep{ritchie2001current,amin2020survey}. Современные языковые модели обладают гораздо более широким покрытием, однако недавние работы показывают, что беглость текста не равна сильной юмористической компетентности. Крупные оценки по-прежнему помещают модельные системы ниже сильных человеческих результатов, а исследования каламбуров и понимания шуток показывают, что модели часто надежнее воспроизводят поверхностную форму юмора, чем его внутренний механизм \citep{Zhang2024HumorAI,xu2024good,cocchieri2025you,loakman2025comparing}.

Иными словами, юмор выявляет знакомую слабость современных моделей: они хорошо продолжают вероятный текст, но хуже выполняют открытый поиск по множеству правдоподобных и неожиданных продолжений. Один из ответов на эту слабость состоит в явном моделировании процесса рассуждения. Современные системы генерации юмора всё чаще используют поэтапное построение подсказок, творческий поиск или декомпозицию на основе теории, а не одношаговое завершение текста. Такие подходы мотивированы простым наблюдением: хорошие шутки часто возникают не через выбор первого вероятного продолжения, а через исследование ассоциаций, переосмыслений, контрастов и форм панчлайна \citep{chen2023prompt,zhong2024let,Tikhonov2024,Kim2025,ZhangJ2025}.

Однако большинство таких систем задаёт структуру рассуждения вручную: исследователь определяет число этапов, тип промежуточных представлений и порядок действий модели. Параллельное направление в постобучении предлагает другую возможность. В верифицируемых доменах обучение с подкреплением использовалось для извлечения развернутого поведения рассуждения напрямую, а не только через ручное построение подсказок или фиксированные процессы \citep{Lightman2023,DeepSeek2025}. Методы без внешнего верификатора, такие как VeriFree и RLPR, дополнительно показывают, что часть этой логики может быть перенесена за пределы математики и программирования, если заменить внешнего верификатора внутренними вероятностными целями \citep{Zhou2025,Yu2025}.

Тем не менее для юмора такие формулировки сразу сталкиваются с несоответствием самой природе задачи. Они часто строятся вокруг предположения об одном эталонном ответе, тогда как юмор естественно является задачей с несколькими эталонными ответами: для запроса вроде ``write a joke about a frog'' существует много приемлемых ответов, и любой набор данных фиксирует только некоторые из них. Данная работа мотивирована именно этим несоответствием.

Центральная идея работы состоит в том, что генерацию юмора следует формулировать не как точную имитацию одной наблюдаемой шутки, а как распределение вероятностной массы по множеству приемлемых шуток при заданном запросе. Этот сдвиг важен и теоретически, и практически. Теоретически он меняет способ определения обучения без внешнего верификатора в творческом домене. Практически он меняет требования к данным: вместо изолированных пар запрос--ответ обучение требует наборов эталонных ответов, связанных с одним запросом и содержащих несколько правдоподобных шуток.

Цель данной работы --- разработать и экспериментально проверить конвейер многореференсного обучения генерации юмора без внешнего верификатора. Для этого решаются следующие задачи:
\begin{enumerate}
    \item построить корпус шуток и очистить его от точных дублей и почти совпадающих текстов;
    \item сформировать слабые запросы на основе ключевых слов и найти для них несколько эталонных шуток;
    \item сформулировать многореференсное обобщение VeriFree для генерации юмора;
    \item заменить сырую вероятность полной последовательности устойчивой нормированной логарифмической массой эталонных ответов;
    \item реализовать обучение LoRA-адаптеров для моделей Qwen3-1.7B и Qwen3-4B;
    \item сгенерировать кандидатные ответы базовых и обученных моделей;
    \item провести автоматическую попарную оценку и проанализировать полученные результаты.
\end{enumerate}

Работа имеет смешанный исследовательско-программный характер. С исследовательской стороны предлагается многореференсная постановка обучения без внешнего верификатора для открытой творческой задачи. С программной стороны реализован полный конвейер: сбор и очистка данных, построение векторных представлений, извлечение ключевых слов, поиск эталонных ответов, обучение, генерация кандидатов, попарная оценка и агрегация таблиц лидеров.

Исходная коллекция содержала приблизительно \(664\,000\) текстов шуток. После удаления дублей и почти совпадающих текстов текущий корпус содержит приблизительно \(570\,000\) шуток. На основе этого корпуса были построены около \(14\,600\) многореференсных запросов. В экспериментальной части обучались LoRA-адаптеры для моделей Qwen3-1.7B и Qwen3-4B. Качество оценивалось автоматически с помощью попарного сравнения ответов моделью Gemini 2.5 Flash и последующей агрегации через долю побед и модель Брэдли--Терри.

Основной эмпирический результат состоит в том, что эффект MRVF зависит от размера базовой модели. Для Qwen3-1.7B результаты оказались смешанными: MRVF-модель с рассуждением заняла второе место, но не превзошла базовую модель с рассуждением. Для Qwen3-4B результаты оказались положительными: обе MRVF-модели превзошли обе базовые модели, а MRVF-модель с рассуждением заняла первое место. Это поддерживает гипотезу о том, что многореференсная логарифмическая цель может давать полезный обучающий сигнал для генерации юмора, если базовая модель достаточно сильна.

Работа устроена следующим образом. В первой главе рассматриваются исследования вычислительного юмора, генерации юмора большими языковыми моделями, оценки юмора, рассуждения и обучения без внешнего верификатора. Во второй главе описывается построение данных, программный конвейер и формальная методология MRVF. В третьей главе приводятся экспериментальная постановка, параметры обучения, протокол генерации и протокол оценки. В четвёртой главе представлены результаты и их интерпретация. В заключении суммируются полученные выводы, ограничения и направления дальнейшей работы.